\documentclass[a4paper, 12pt]
{article}
\usepackage{amsmath, amssymb, esvect, mathtools, enumerate, enumitem}

\newcommand{\inv}[1]{#1^\times}
\newcommand{\z}{\mathbb{Z}}
\newcommand{\rpol}{R[x]}
\newcommand{\zpol}{\mathbb{Z}[x]}
\newcommand{\pol}[2]{#1[#2]}
\newcommand{\nonzeroset}[1]{#1 \backslash \{0\}}
\newcommand{\setwithzero}[1]{#1_{0}}
\newcommand{\letter}[1]{$#1$}
\newcommand{\seq}[2]{#1_{1}, #1_{2}, \ldots #1_{#2}}


\begin{document}

	\textbf{Math 322: Commutative Algebra}
	
	\hfill


	\textbf{Chapter 1: Euclidean Domains and Principal Ideal Domains}

	\hfill
	
	\begin{itemize}
		\item \textbf{Ring} Abelian group $(R,+)$ equipped with a multiplication that satisfies left and right distributivity property and contains a  multiplicative identity
		
		\item \textbf{Commutative Ring} If the multiplication is commutative
		
		\item \textbf{Zero Divisor} Let R be a commutative ring. An element $a \in R \backslash \{0\}$ is a non-zero divisor if there exists $b \in R \backslash \{0\}$ such that $ab=0.$
		
		\item \textbf{Integral Domain} A commutative ring with no zero divisors, written ID.  
		
		\item\textit{ Property of ID} Multiplicative Cancellation:
		If $0 \not = a \in R$ and $b,c \in R$ such that $ab=ac$ then $b=c.$
		
		\item \textbf{Ideal} A subgroup $I$ of $(R,+)$ such that $ab, ba \in I$ for any $a \in R$ and $b \in I. $
		
		\item \textbf{Generating Set} A generating set for an ideal $I$ is a subset $S$ of $I$, such that the elements of $S$ $R$-generate $S.$
		
		\item \textbf{Ideal of $R$ generated by $S$} If $S=\{s_1, s_2, \cdots s_n\}$ is finite then $I=s_1R + s_2R + \cdots s_nR$ is the ideal of $R$ generated by $S.$
		
		\item \textbf{Principal} An ideal $I$ is principal if $I$ can be generated by single element.
		
		\item Given two ideals $I,J$ in a commutative ring $R,$ we can construct new ideals: 
		$$I \cap J,$$ 
		$$I+J, $$
		$$IJ=\left\{\sum_\text{finite}ab: a\in I, b\in J\right\}$$
		
		\item\textit{ Notation.} We write $a \equiv b \mod I$ for $a-b \in I.$
		
		\item \textbf{Principal Ideal Domain} An integral domain such that any ideal is principal.
		
		\item Main instances of PID's:
		\begin{enumerate}
			\item Ring of Integers
			\item Polynomial rings in one variable over a field
		\end{enumerate}
		
		\item \textbf{Euclidean Function} Let \letter{R} be an integral domain. A Euclidean Function on \letter{R} is a function $v: R \backslash \{0\} \rightarrow \mathbb{N}_0$ such that the following hold:
					\begin{enumerate}[label=(\roman*)]
						\item $v(ab) \geq v(a)$ for all $a,b\in R \backslash \{0\}$
						\item for any $a\in R$ and any $b\in R \backslash \{0\},$ there exists $q,r\in R$ such that $a=qb+r$ and either $r=0$ or $v(r)<v(b)$
					\end{enumerate}
		\item \textbf{Euclidean Domain} ID \letter{R} such that there is a Euclidean Function on \letter{R.}
		
		\item \textit{Important ED:}
		Ring of Gaussian integers $\pol{\z}{i}=\{a+bi: a,b\in \mathbb{Z}\}$ paired with the Euclidean Function 
		$$v: \mathbb{Z}[i] \backslash \{0\} \longrightarrow \mathbb{N}_0, v(z) = |z|^2  \quad \text{for all} \quad z \in \pol{\z}{i}$$
		
		\item ED $\implies$ PID
		
		\item \textit{Important theorem relating a ring and its polynomial ring}
		Let \letter{R} be a ring and $\rpol$ the polynomial ring over \letter{R}. TFAE
					\begin{enumerate}[label=(\roman*)]
						\item $\rpol$ is a ED
						\item $\rpol$ is a PID
						\item $R$ is a field
					\end{enumerate}
					
		\item \textit{Non-example to above theorem}
		Consider the ring $\zpol$ of polynomials with integer coeffiecients. Recall that $\mathbb{Z}$ is not a field. Show $\zpol$ is not a ED by showing second axiom of ED fails using the functions $f=x^3-3x+1$ and $g=2x.$

		
		In the next three definitions, let $R$ be a commutative unital ring, with multiplicative identity 1 and additive identity 0.
		
		\item \textbf{Inveritible} An element $a\in R$ is invertible if there exists $b\in R$ such that $ab=1.$
		
		 $\inv{R}$ is the set of all invertible elements of $R$ and $a^{-1}$ is the inverse of $a.$
		
		\item \textbf{Divides/Factor} If $a,b \in R$ are both non-zero, we say $a$ divides $b$  $a$ is a factor of  $b$ if there exists $c\in R$ such that $b=ac.$ Write $a|b.$
		
		\item  \textbf{Associated} Two non-zero elements $a,b\in R$ are associated if only if $a | b$ \textbf{and} $b|a.$ Write $a \sim b.$
		
		\item \textit{Useful divisibility proposition}
		Let $R$ be a commutative ring and $a,b\in R\backslash \{0\}.$ Then $$a|b \iff bR \subseteq aR.$$
		
		When considering questions of divisibility, we will generally work with IDs only.
		
		\item \textit{Remarks}
		\begin{itemize}
			\item We can also refer to an invertible element as a \textbf{unit} in a ring.
			\item $u \in\inv{R} \iff uR=R$
			\item $\inv{R}$ forms a mulitplicative group
			\item The relation `$a$ is associated to $b$ in $R$` is an equivalence relation
		\end{itemize}
		
		\item \textit{Characterisation of invertible elements of a polynomial IDs} Let \letter{R} an ID. Then $\inv{\rpol}=\inv{R}$ (All	 the invertible elements in $\rpol$ are the constant polynomials.
		
		\item \textit{Characterisation of invertible elements of a ED}
		Let $R$ be a ED with Euclidean function $v.$ The following hold
		\begin{enumerate}[label=(\roman*)]
			\item $v(1) =\{v(a): a \in \nonzeroset{R}\}$
			\item $a\in R$ is invertible if and only if $v(a)=v(1).$
		\end{enumerate}
		
		\item Every field has an infinite (countably) many Euclidean Funcitons.
		
		For any field \letter{R}, show the constant function 
		$$v_n: \nonzeroset{R} \rightarrow \setwithzero{\mathbb{N}}, \quad v_n(a)=n, \quad \forall a\in \nonzeroset{R}$$
		is a Euclidean function.
		
		\item \textbf{Common Divisor/Greatest Common Divisor}
		
		Let \letter{R} be a commutative ring and $\seq{a}{n} \in R \quad (n \geq 2.)$
		
		An element $a \in R$ is a common divisor of $\seq{a}{n}$ if $d | a_i$ for all $i.$ It is a greatest common divisor if it is a common divisor of $\seq{a}{n}$ and any common divisor $c$ is also a divisor of $d$ 
		
		\item \textit{Nice result about GCDs} 
		
		Let \letter{R} be an ID and $\seq{a}{n} \in R.$ Suppose that $\gcd(\seq{a}{n})$ exists in \letter{R}, say $d =\gcd(\seq{a}{n}) \in R.$
		
	\begin{enumerate}[label=(\roman*)]
		\item  Any gcd of $\seq{a}{n}$ is associated to \letter{d}
		\item \letter{d} divides any $R-$linear combination of $\seq{a}{n}.$ If an $R-$linear combination od $\seq{a}{n}$ is a common divisior of $\seq{a}{n}$ then it is a gcd of $\seq{a}{n}.$
	\end{enumerate}
	
	\item \textit{Euclidean Algorithm} Let \letter{R} be a Euclidean domain with Euclidean function \letter{v}. Let $a,b\in R$ and suppose that $v(a) \geq v(b).$ By definition of Euclidean Domain, there exists $q_1, r_1 \in R$ such that 
	$$a=q_1b+r_1, \quad \text{and} \quad r_1=0 \quad \text{or} \quad v(r_1)<v(b)$$
	
	If $r_1=0$ then $b|a$ and so $b=\gcd(a,b).$ Otherwise repeat the above process defining $q_2, r_2$ similarly unitl we eventually reach $r_n=0.$
	\end{itemize}

	
	
\end{document}
